% ============================================================================
% resumen_ejecutivo.tex -- Executive Summary
% Global-E-Shop Churn Prediction System
% Universidad de La Sabana . Ciencia de Datos II . 2026.1
% Docente: Sergio Amortegui
%
% Compile: pdflatex resumen_ejecutivo.tex
% ============================================================================
\documentclass[11pt,a4paper]{article}

\usepackage[utf8]{inputenc}
\usepackage[T1]{fontenc}
\usepackage[english]{babel}
\addto\captionsenglish{%
  \renewcommand{\contentsname}{Tabla de Contenidos}%
  \renewcommand{\figurename}{Figura}%
  \renewcommand{\tablename}{Tabla}%
}

\usepackage{geometry}
\geometry{top=2.5cm, bottom=2.5cm, left=2.8cm, right=2.8cm, headheight=15pt}

\usepackage{xcolor}
\definecolor{primary}{RGB}{0,84,166}
\definecolor{accent}{RGB}{200,40,50}
\definecolor{success}{RGB}{30,140,120}
\definecolor{darkgray}{RGB}{60,65,70}
\definecolor{lightblue}{RGB}{232,242,255}

\usepackage{tcolorbox}
\usepackage{booktabs}
\usepackage{array}
\usepackage{tabularx}
\usepackage{parskip}
\usepackage{titlesec}
\usepackage{fancyhdr}
\usepackage{url}
\usepackage{amsmath}
\usepackage{amssymb}
\usepackage{enumitem}
\usepackage{graphicx}
\usepackage{ragged2e}

% ── Encabezado / pie ────────────────────────────────────────────────────────
\pagestyle{fancy}
\fancyhf{}
\fancyhead[L]{\small\color{darkgray} Global-E-Shop $|$ Churn Prediction System}
\fancyhead[R]{\small\color{darkgray} Resumen Ejecutivo $\cdot$ 2026.1}
\fancyfoot[C]{\small\color{darkgray}\thepage}
\renewcommand{\headrulewidth}{0.4pt}

% ── Secciones ───────────────────────────────────────────────────────────────
\titleformat{\section}{\large\bfseries\color{primary}}{}{0em}{}[\color{primary}\titlerule]
\titleformat{\subsection}{\normalsize\bfseries\color{darkgray}}{}{0em}{}
\titlespacing{\section}{0pt}{14pt}{6pt}
\titlespacing{\subsection}{0pt}{10pt}{4pt}

% ── tcolorbox simples (sin skins) ────────────────────────────────────────────
\newcommand{\kpibox}[2]{%
  \begin{tcolorbox}[colback=lightblue,colframe=primary,
    title={\bfseries\color{primary}#1},
    left=6pt, right=6pt, top=5pt, bottom=5pt, boxrule=0.8pt, arc=4pt]
    #2
  \end{tcolorbox}%
}
\newcommand{\alertbox}[2]{%
  \begin{tcolorbox}[colback=red!7!white,colframe=accent,
    title={\bfseries\color{accent}#1},
    left=6pt, right=6pt, top=4pt, bottom=4pt, boxrule=0.8pt, arc=4pt]
    #2
  \end{tcolorbox}%
}
\newcommand{\successbox}[2]{%
  \begin{tcolorbox}[colback=green!8!white,colframe=success,
    title={\bfseries\color{success}#1},
    left=6pt, right=6pt, top=4pt, bottom=4pt, boxrule=0.8pt, arc=4pt]
    #2
  \end{tcolorbox}%
}

% ── Columna de tabla ancha ──────────────────────────────────────────────────
\newcolumntype{L}[1]{>{\RaggedRight\arraybackslash}p{#1}}

% ============================================================================
\begin{document}

% ── PORTADA ─────────────────────────────────────────────────────────────────
\thispagestyle{empty}

\begin{center}
  \vspace*{0.5cm}
  {\color{primary}\rule{\linewidth}{2pt}}\\[0.7cm]
  {\LARGE\bfseries\color{primary} Sistema de Prediccion de Churn}\\[6pt]
  {\LARGE\bfseries\color{primary} Global-E-Shop}\\[0.5cm]
  {\large\bfseries Resumen Ejecutivo}\\[0.4cm]
  {\large\color{darkgray} Proyecto de Ciencia de Datos II}\\
  {\color{darkgray} Universidad de La Sabana $\cdot$ Facultad de Ingenieria $\cdot$ 2026.1}\\[4pt]
  {\color{darkgray} Docente: Sergio Amortegui}\\[0.5cm]
  {\color{primary}\rule{\linewidth}{2pt}}
\end{center}

\vspace{0.5cm}

% ── PROPÓSITO ───────────────────────────────────────────────────────────────
\kpibox{Proposito de este Documento}{%
  Este resumen esta dirigido al \textbf{Gerente de Marketing} de Global-E-Shop.
  Presenta la decision de negocio que sustenta el sistema, los resultados
  validados del modelo de prediccion de abandono (\textit{churn}), el impacto
  economico estimado y los proximos pasos recomendados.
  Los detalles tecnicos completos se encuentran en el \textit{Informe Tecnico APA}
  (24 pags.) y en el \textit{Analytics Contract} (15 pags.).
}

\vspace{0.4cm}

% ── 1. PROBLEMA DE NEGOCIO ───────────────────────────────────────────────────
\section{Problema de Negocio y Decision}

Global-E-Shop enfrenta una tasa de abandono de clientes del $\approx 10\%$ trimestral.
Adquirir un cliente nuevo cuesta entre 5 y 7 veces mas que retener uno existente
(Reichheld \& Schefter, 2000). Sin embargo, intervenir indiscriminadamente
sobre toda la base de clientes genera campanas costosas con ROI negativo.

\textbf{Decision formal:} Para cada cliente $c$ al cierre del mes, el sistema elige
una accion $a \in \{a_0\text{ (no actuar)},\ a_1\text{ (email)},\ a_2\text{ (descuento 10\%)},
\ a_3\text{ (descuento 20\%)}\}$, maximizando:
\[
  U(a,\omega) \;=\; \mathrm{CLV} \cdot \mathbf{1}[\text{retenido} \mid \omega, a]
                 \;-\; \mathrm{costo}(a)
\]
donde $\omega \in \{\text{Churn, No-Churn}\}$ y $\mathrm{CLV} \approx \$480$
por cliente activo (valor de ciclo de vida promedio).

% ── 2. RESULTADOS DEL MODELO ─────────────────────────────────────────────────
\section{Resultados Validados del Modelo}

\kpibox{KPIs Clave --- XGBoost con SMOTE (Test Set, B=1{,}000 Bootstrap)}{%
  \vspace{2pt}
  \renewcommand{\arraystretch}{1.4}
  \begin{tabular}{lcc}
    \toprule
    \textbf{Metrica} & \textbf{Valor} & \textbf{Umbral Minimo} \\
    \midrule
    \textbf{Recall (Churn)} & \textbf{0.82} [IC 95\%: 0.80,\,0.84] & $\geq 0.80$ $\checkmark$ \\
    Precision               & 0.45                                  & $\geq 0.30$ $\checkmark$ \\
    AUC-ROC                 & 0.87                                  & $\geq 0.65$ $\checkmark$ \\
    AUC-PR                  & 0.52                                  & ---         \\
    \textbf{Lift@20\%}      & \textbf{2.4x}                         & $\geq 1.50\times$ $\checkmark$ \\
    \bottomrule
  \end{tabular}
}

\vspace{0.4cm}

\textbf{Comparativa contra baselines honestos:}

\vspace{0.3cm}
\renewcommand{\arraystretch}{1.35}
\begin{tabularx}{\linewidth}{L{4.2cm}cccc}
  \toprule
  \textbf{Metodo} & \textbf{Recall} & \textbf{Precision} & \textbf{AUC-ROC} & \textbf{Lift@20\%} \\
  \midrule
  Baseline RFM (p75)      & 0.45 & 0.11 & 0.62 & 1.20x \\
  Recency Fija (60 dias)  & 0.38 & 0.09 & 0.58 & 1.10x \\
  Frecuencia Baja ($<$2)  & 0.41 & 0.10 & 0.60 & 1.15x \\
  RFM Compuesto           & 0.47 & 0.12 & 0.63 & 1.22x \\
  \midrule
  \textbf{XGBoost (sistema)} & \textbf{0.82} & \textbf{0.45} & \textbf{0.87} & \textbf{2.4x} \\
  \bottomrule
\end{tabularx}

\vspace{0.4cm}

El modelo supera al mejor baseline en \textbf{+0.35 puntos de Recall} y
\textbf{+1.18x de Lift}. Todos los resultados incluyen Intervalos de Confianza
Bootstrap al 95\% (B=1{,}000 iteraciones, alpha=0.05).

% ── 3. IMPACTO ECONÓMICO ─────────────────────────────────────────────────────
\section{Impacto Economico Estimado}

Con una base de \textbf{5{,}000 clientes activos} y tasa de churn del 10\%
($\approx 500$ clientes en riesgo por trimestre):

\vspace{0.3cm}
\renewcommand{\arraystretch}{1.4}
\begin{tabularx}{\linewidth}{L{5.4cm}cc}
  \toprule
  \textbf{Escenario} & \textbf{Churners recuperados} & \textbf{Ingresos retenidos} \\
  \midrule
  Sin modelo (campana masiva)           & $\sim$190 (Recall$\approx$0.38) & \$91{,}200  \\
  Baseline RFM                          & $\sim$225 (Recall$\approx$0.45) & \$108{,}000 \\
  \textbf{XGBoost (sistema actual)}     & $\sim$\textbf{410} (Recall$\approx$0.82) & \textbf{\$196{,}800} \\
  \bottomrule
\end{tabularx}

\vspace{0.4cm}

\successbox{Beneficio Neto Estimado vs.\ Baseline RFM}{%
  \textbf{+\$88{,}800 en ingresos retenidos} por trimestre al pasar del baseline RFM al
  sistema XGBoost, asumiendo CLV = \$480 y costo promedio de campana = \$12/cliente.\\[4pt]
  \textit{Supuesto: eficacia del 40\% de la accion sobre churners correctamente
  identificados. Calculo ilustrativo; se recomienda A/B test para validacion en produccion.}
}

% ── 4. ARQUITECTURA DEL SISTEMA ───────────────────────────────────────────────
\section{Arquitectura del Sistema}

\begin{center}
\small
\begin{tabular}{ccccccc}
  \texttt{Seed} & $\rightarrow$ & \texttt{ETL} & $\rightarrow$ & \texttt{EDA/QA} & $\rightarrow$ & \texttt{Features} \\[4pt]
  & & & & & & $\downarrow$ \\[4pt]
  \texttt{API FastAPI} & $\leftarrow$ & \texttt{Modelo} & $\leftarrow$ & \texttt{Evaluacion} & $\leftarrow$ & \texttt{Train/SMOTE} \\
\end{tabular}
\end{center}

\vspace{0.3cm}
Tres garantias de produccion:

\begin{enumerate}[leftmargin=1.5em, itemsep=3pt]
  \item \textbf{Zero Data Leakage}: features sobre $W_\text{obs} = [t_0 - 90\text{d},\ t_0]$;
        etiquetas sobre $W_\text{pred} = (t_0,\ t_0 + 30\text{d}]$. Sin solapamiento temporal.
  \item \textbf{Reproducibilidad}: semilla global \texttt{245573}, dependencias fijadas en
        \texttt{requirements.txt}, artefactos con hash MD5 en \texttt{artifacts/manifest.json}.
  \item \textbf{Trazabilidad}: cada prediccion queda registrada en \texttt{ml\_predictions\_log}
        con version del modelo, timestamp y score de probabilidad.
\end{enumerate}

% ── 5. PRÓXIMOS PASOS ────────────────────────────────────────────────────────
\section{Proximos Pasos Recomendados}

\renewcommand{\arraystretch}{1.4}
\begin{tabularx}{\linewidth}{L{0.5cm}L{4.0cm}L{8.2cm}}
  \toprule
  \textbf{P} & \textbf{Accion} & \textbf{Descripcion} \\
  \midrule
  1 & \textbf{A/B Test en produccion}
    & Dividir la proxima cohorte en riesgo: 50\% guiada por el modelo,
      50\% por regla RFM. Medir Recall real y ROI a 60 dias. \\
  2 & \textbf{Monitoreo de drift}
    & PSI (\textit{Population Stability Index}) mensual sobre Recency y Frequency.
      Reentrenar si PSI $> 0.25$ (Gama et al., 2014). \\
  3 & \textbf{Enriquecimiento de senales}
    & Integrar logs de navegacion web (abandono de carrito, sesiones sin compra)
      como features adicionales. \\
  4 & \textbf{Despliegue en nube}
    & Containerizar la API FastAPI (Docker) y desplegar en Cloud Run o ECS.
      Activar cache de predicciones para clientes recurrentes. \\
  \bottomrule
\end{tabularx}

% ── 6. GARANTÍAS DE CALIDAD ───────────────────────────────────────────────────
\section{Garantias de Calidad de Datos (Stop/Go)}

\alertbox{Criterios Formales Verificados Antes de Liberar el Modelo}{%
  \begin{enumerate}[leftmargin=1.5em, itemsep=2pt]
    \item \textbf{Integridad de llaves}: \textit{dup\_rate} = 0\% en todas las tablas
          (\texttt{v\_key\_integrity}).
    \item \textbf{Integridad referencial}: \textit{orphan\_rate} = 0\% en todos los joins
          (\texttt{v\_key\_integrity}).
    \item \textbf{Cobertura del Feature Store}: 100\% de clientes con etiqueta tienen
          features calculadas (\texttt{v\_feature\_coverage}).
    \item \textbf{Nulos bajo umbral}: ninguna columna supera el 20\% de nulos
          (\texttt{config.yaml: etl.null\_threshold}).
  \end{enumerate}
  \textbf{Veredicto}: \textcolor{success}{\textbf{PASS en los 4 criterios.
  Pipeline habilitado para entrenamiento.}}
}

% ── PIE DE PÁGINA ─────────────────────────────────────────────────────────────
\vfill
\noindent\rule{\linewidth}{0.4pt}\\[4pt]
\begin{minipage}{0.58\linewidth}
  \small\color{darkgray}
  \textbf{Autores:} Juan Manuel Prieto Corredor \&\ Slendy Marieth Grisales Rueda\\
  Facultad de Ingenieria, Universidad de La Sabana\\
  \texttt{juanmanuelprietocorredor@gmail.com}
\end{minipage}
\hfill
\begin{minipage}{0.38\linewidth}
  \RaggedLeft\small\color{darkgray}
  Repositorio:\\
  \url{github.com/juanmanuelpriet/}\\
  \url{Proyecto-ciencia-ade-datos-2}
\end{minipage}

\end{document}
